\chapter{Foundations: the ideas behind object detection}

This chapter builds the vocabulary. If you already know what IoU, anchors, NMS, feature
pyramids, and segmentation are, you can skim. Otherwise, read it once and the code
chapters will read much more easily.

\section{What ``object detection'' actually asks for}
Image \emph{classification} answers ``what is in this picture?'' with a single label.
\textbf{Object detection} is harder: it must answer ``\emph{what} objects are here, and
\emph{where} is each one?'' --- producing, for every object, a \textbf{bounding box}
(a rectangle) and a \textbf{class label}. This codebase then adds two more questions per
object: ``what is its exact outline?'' and ``how far away is it?''.

\section{Bounding boxes and coordinate conventions}
A box is four numbers. But there are several ways to write those four numbers, and mixing
them up is the \#1 bug in detection code.

\begin{center}
\renewcommand{\arraystretch}{1.25}
\begin{tabular}{l l}
\toprule
\textbf{Convention} & \textbf{The four numbers} \\
\midrule
\code{xyxy} & left, top, right, bottom (two corners), x first \\
\code{yxyx} & top, left, bottom, right (two corners), y first \\
\code{xywh} & center-x, center-y, width, height \\
normalized & each value divided by image width/height, so in $[0,1]$ \\
pixels      & raw pixel coordinates \\
\bottomrule
\end{tabular}
\end{center}

\begin{pitfall}
This repo deliberately uses \textbf{two} conventions and converts between them at a known
seam. Ground-truth boxes coming out of the dataset decoders and parsers are
\code{yxyx}, \textbf{normalized} to $[0,1]$. The loss and the label-assignment code
convert to \code{xyxy} in \textbf{pixels}. Whenever you touch box math, first ask
``which convention am I in?''. The codebase is careful to label this everywhere; you
should be too.
\end{pitfall}

\section{Intersection over Union (IoU)}
To say whether a predicted box is ``correct'', we need to measure how well it overlaps a
ground-truth box. The standard measure is \textbf{Intersection over Union}: the area of
overlap divided by the area of their union.

\begin{center}
\begin{tikzpicture}[scale=0.9]
  \fill[bluebg,draw=accent,thick] (0,0) rectangle (2.6,1.8);
  \fill[greenbg,draw=green,thick,opacity=0.7] (1.4,0.6) rectangle (3.8,2.4);
  \fill[amber,opacity=0.35] (1.4,0.6) rectangle (2.6,1.8);
  \node[font=\scriptsize] at (0.7,1.5) {prediction};
  \node[font=\scriptsize] at (3.2,2.1) {ground truth};
  \node[font=\scriptsize,text=amber] at (2.0,1.2) {overlap};
  \node[font=\small] at (6.6,1.2) {$\text{IoU}=\dfrac{\text{area of overlap}}{\text{area of union}}\in[0,1]$};
\end{tikzpicture}
\end{center}

\noindent IoU is 1 for a perfect match and 0 for no overlap. A common rule is ``a
detection counts as correct if its IoU with a true box is at least 0.5''. IoU appears
everywhere: in the metric that scores the model, in NMS, and --- in a refined form called
\textbf{CIoU} --- inside the box loss.

\section{Confidence scores and NMS}
Each detection comes with a \textbf{confidence} (or score) in $[0,1]$: how sure the model
is. A detector fires on the same object from several adjacent grid cells, so we get a
cluster of overlapping boxes. \textbf{Non-maximum suppression} cleans this up: sort by
score, keep the top box, delete every other box that overlaps it too much (IoU above a
threshold), repeat.

\begin{hood}
Here NMS runs in \file{models/detection\_generator.py} with \code{max\_boxes=300},
\code{nms\_thresh=0.65}, \code{score\_thresh=0.05}, and it is \textbf{per-class}: classes
do not suppress each other, so an overlapping ``person'' and ``bag'' both survive.
\end{hood}

\section{The grid idea: how YOLO predicts}
YOLO (``You Only Look Once'') does not propose-then-classify. It overlays a \textbf{grid}
on the feature map and asks \emph{every cell} to predict whether an object is centered
there, and if so its box, class, and (here) polygon and distance. One forward pass yields
all predictions --- fast and simple.

\begin{teacher}
\textbf{Stride.} A backbone shrinks the image. If a feature map is 8$\times$ smaller than
the input, we say it has \textbf{stride 8}: each cell ``sees'' an $8\times8$ pixel patch.
Small strides (fine grids) catch small objects; large strides (coarse grids) catch large
ones.
\end{teacher}

\section{Feature pyramids and multiple scales}
A single grid can't handle a tiny sign and a huge truck equally well. So detectors use a
\textbf{feature pyramid}: several feature maps at different strides. This codebase uses
three, at strides \textbf{8, 16, and 32} (grids of $84\times84$, $42\times42$,
$21\times21$ for a 672-pixel input). Every head predicts at all three. The decoder that
builds and fuses these maps is an \textbf{FPN-PAN} --- covered in the architecture chapter.

\section{Anchor-free detection}
Older detectors placed several preset box shapes (``anchors'') at each cell and predicted
offsets from them. YOLOv8 is \textbf{anchor-free}: each cell has a single \emph{anchor
point} at its center, and the box head predicts the distances from that point to the four
sides. The anchor points are simply the cell centers, $(i+0.5)\cdot\text{stride}$,
built inside the loss and the detection generator.

\section{Two ways to describe an object's shape}
A bounding box is a crude outline. \textbf{Segmentation} describes the actual shape. There
are two common encodings:

\begin{itemize}[leftmargin=1.4em,itemsep=2pt]
  \item a \textbf{mask} --- a per-pixel ``inside/outside'' map (heavy, but exact);
  \item a \textbf{polygon} --- an ordered list of boundary points (compact).
\end{itemize}

\noindent This codebase uses a special \emph{radial} polygon format from
``PolyYOLO'': 24 vertices arranged at fixed $15^\circ$ angles around the box center. It
is light enough to predict with three small heads and is the subject of
Chapter~\ref{ch:polygons}.

\section{Estimating distance}
The last extra task is \textbf{monocular distance estimation}: from a single RGB image,
predict how far each object is, in meters. This is learned from a separate labeled dataset
and predicted by a one-channel head on a log scale. It is trained jointly with detection,
merged at the batch level --- a neat trick we will see in the data-pipeline chapter.

\section{How a model ``learns where things are''}
One subtlety underlies all of detection training. The model makes predictions at thousands
of grid cells, but a training image has only a handful of objects. Which cells should be
held responsible for which object? This is \textbf{label assignment}, and modern YOLO uses
a smart, learnable scheme called \textbf{Task-Aligned Assignment (TAL)} that picks, for
each object, the cells whose predictions are already both well-localized and
well-classified. We devote a section to it in the losses chapter --- it is the heart of how
the model is trained.

\medskip
\noindent That is the whole vocabulary. Armed with it, let's open the codebase.
